\section{Step 65 Distinguish a direct length from a repeat operator}\label{png-huff-s65}
Implement \texttt{huff\_direct\_length} after the supplied code-length symbol call. Values zero through fifteen provide one direct length. Retain that value, set repetition count one and select \texttt{store\_lengths}. Larger decoded values select the repeat dispatch instead.

The first decoded symbol in the dynamic-literals fixture is one, meaning that literal/length symbol zero will have a one-bit code. At this checkpoint the combined cursor has not advanced. Reusing the data-literal emitter here would incorrectly write a row byte instead of a code length.

\paragraph{Observe} Use \texttt{s65-1}; stop before \texttt{store\_lengths}. Inspect the state described above before continuing.

\Needspace{7\baselineskip}
\paragraph{Check} What does decoded symbol one mean in this table?
\begin{itemize}
\item[A.] A literal data byte one.
\item[B.] One code length whose value is one.
\item[C.] A distance of one.
\end{itemize}
